% !TEX root = ../../ctfp-print.tex

\lettrine[lhang=0.17]{型}{を組み合わせる2つの基本的な}方法を見てきた：直積と余積を使う方法だ。
日常のプログラミングにおける多くのデータ構造は、この2つの仕組みだけを使って構築できることがわかっている。
この事実は実践的に重要な意味を持つ。データ構造の多くの性質は合成可能だ。例えば、基本型の値を
等値比較する方法を知っていて、その比較を直積型と余積型に一般化する方法を知っていれば、合成型の
等値演算子の導出を自動化できる。Haskellでは、合成型の大きな部分集合に対して、等値比較、順序比較、
文字列との相互変換などを自動的に導出できる。

プログラミングにおける直積型と和型をより詳しく見ていこう。

\section{直積型}

プログラミング言語における2つの型の直積の標準的な実装はペアだ。Haskellではペアはプリミティブな
型コンストラクタだ；C++では標準ライブラリで定義された比較的複雑なテンプレートだ。

\begin{figure}[H]
  \centering
  \includegraphics[width=0.35\textwidth]{images/pair.jpg}
\end{figure}

\noindent
ペアは厳密には可換でない：ペア\code{(Int, Bool)}は同じ情報を持っていても、ペア
\code{(Bool, Int)}の代わりには使えない。しかし、同型を除いては可換だ――同型は
\code{swap}関数（これ自身が逆関数）によって与えられる：

\src{snippet01}
2つのペアは、同じデータを異なる形式で格納しているだけと考えることができる。ビッグエンディアン対
リトルエンディアンのようなものだ。

ペアをネストすることで、任意の数の型を直積に組み合わせることができるが、より簡単な方法がある：
ネストされたペアはタプルと等価だ。これは、ペアをネストする異なる方法が同型であるという事実の帰結だ。
例えば、3つの型\code{a}、\code{b}、\code{c}をこの順に直積に組み合わせるには、2通りの方法がある：

\src{snippet02}

または
\src{snippet03}
これらの型は異なる――一方を期待する関数に他方を渡すことはできない――しかし、それらの要素は一対一
対応している。一方を他方に写す関数がある：

\src{snippet04}
そしてこの関数は可逆だ：

\src{snippet05}
よって、これは同型だ。これらは同じデータを再パッケージする異なる方法に過ぎない。

直積型の生成を型上の二項演算として解釈できる。その観点からすると、上記の同型はモノイドで見た
結合律によく似ている：
\[(a * b) * c = a * (b * c)\]
ただし、モノイドの場合は直積の2つの合成方法が等しかったのに対し、ここでは「同型を除いて」のみ
等しい。

同型で満足でき、厳密な等値を主張しないなら、さらに進んで、ユニット型\code{()}が
1が乗算の単位元であるのと同様に直積の単位元であることを示せる。
確かに、ある型\code{a}の値をユニットとペアにしても情報は追加されない。型：

\src{snippet06}
は\code{a}と同型だ。以下がその同型だ：

\src{snippet07}

\src{snippet08}
これらの観察は、$\Set$（集合の圏）が\newterm{モノイド圏}であると言うことで形式化できる。
これは圏であると同時にモノイドでもある圏で、対象を掛け合わせることができる（ここでは
デカルト積をとる）。モノイド圏についてはまた詳しく話し、完全な定義を後で述べる。

Haskellでは、特に和型と組み合わせる場合に、より一般的な方法で直積型を定義できる。
それは複数の引数を持つ名前付きコンストラクタを使う方法だ。例えば、ペアは次のように
別の方法で定義できる：

\src{snippet09}
ここで\code{Pair a b}は他の2つの型\code{a}と\code{b}によってパラメータ化された型の名前で；
\code{P}はデータコンストラクタの名前だ。2つの型を\code{Pair}型コンストラクタに渡すことで
ペア型を定義する。適切な型の2つの値をコンストラクタ\code{P}に渡すことでペアの値を構築する。
例えば、\code{String}と\code{Bool}のペアとして値\code{stmt}を定義してみよう：

\src{snippet10}
最初の行は型シグネチャだ。型コンストラクタ\code{Pair}を使い、\code{Pair}のジェネリックな
定義における\code{a}と\code{b}を\code{String}と\code{Bool}で置き換えている。2行目は
具体的な文字列と具体的なブール値をデータコンストラクタ\code{P}に渡すことで実際の値を定義する。
型コンストラクタは型を構築するために使われ；データコンストラクタは値を構築するために使われる。

Haskellでは型コンストラクタとデータコンストラクタの名前空間は分かれているため、次のように
同じ名前を両方に使うのをよく見かける：

\src{snippet11}
十分に目を細めれば、組み込みのペア型をこの種の宣言の変形と見なすこともできる。ただし、名前
\code{Pair}が二項演算子\code{(,)}に置き換えられている。実際、\code{(,)}を他の名前付き
コンストラクタと同様に使い、前置記法でペアを作れる：

\src{snippet12}
同様に、\code{(,,)}を使って3つ組を作ることもできる。

汎用のペアやタプルを使う代わりに、次のように特定の名前付き直積型を定義することもできる：

\src{snippet13}
これは\code{String}と\code{Bool}の直積に過ぎないが、独自の名前とコンストラクタが与えられている。
この宣言スタイルの利点は、同じ内容を持つが異なる意味と機能を持ち、互いに代替できない型を多数定義
できることだ。

タプルと複数引数のコンストラクタでプログラミングすると、どのコンポーネントが何を表すかを追跡するのが
面倒でエラーが起きやすくなる。コンポーネントに名前をつける方が好ましいことが多い。名前付きフィールドを
持つ直積型はHaskellでは\newterm{レコード}と呼ばれ、Cでは\code{struct}と呼ばれる。

\section{レコード}

簡単な例を見てみよう。化学元素を、2つの文字列（名前と元素記号）と整数（原子番号）を1つの
データ構造に組み合わせて記述したい。タプル\code{(String, String, Int)}を使い、どのコンポーネントが
何を表すかを覚えておくことができる。コンポーネントはパターンマッチングで抽出する。例えば、
元素の記号がその名前のプレフィックスであるかどうかをチェックするこの関数のように
（\textbf{He}が\textbf{Helium}のプレフィックスであるように）：

\src{snippet14}
このコードはエラーが起きやすく、読みにくく保守しにくい。レコードを定義する方がずっと良い：

\src{snippet15}
2つの表現は同型であり、それはこの2つの変換関数（互いの逆関数）によって示される：

\src{snippet16}

\src{snippet17}
レコードフィールドの名前はそのフィールドにアクセスするための関数としても機能することに注目しよう。
例えば、\code{atomicNumber e}は\code{e}から\code{atomicNumber}フィールドを取得する。
\code{atomicNumber}を次の型の関数として使う：

\src{snippet18}
\code{Element}のレコード構文を使うと、関数\code{startsWithSymbol}はより読みやすくなる：

\src{snippet19}
Haskellのトリックとして\code{isPrefixOf}関数をバッククォートで囲んで中置演算子に変えることで、
ほぼ文章のように読めるようにすることもできる：

\src{snippet20}
この場合、中置演算子は関数呼び出しより優先順位が低いため、括弧は省略できる。

\section{和型}

集合の圏における直積が直積型を生み出すのと同様に、余積は和型を生み出す。Haskellにおける
和型の標準的な実装は次のとおりだ：

\src{snippet21}
ペアと同様に、\code{Either}は可換（同型を除いて）で、ネスト可能であり、ネストの順序は
（同型を除いて）無関係だ。よって、例えば3つ組の和型相当を定義できる：

\src{snippet22}
などなど。

$\Set$は余積に関しても（対称）モノイド圏であることがわかる。二項演算の役割を直和が果たし、
単位元の役割を始対象が果たす。型の言葉では、\code{Either}がモノイド演算子で、\code{Void}
（空の型）が中立元だ。\code{Either}を加算、\code{Void}をゼロと考えることができる。
確かに、和型に\code{Void}を加えてもその内容は変わらない。例えば：

\src{snippet23}
は\code{a}と同型だ。なぜなら、この型の\code{Right}版を構築する方法がないからだ――
\code{Void}型の値は存在しない。\code{Either a Void}の唯一の要素は\code{Left}コンストラクタを
使って構築され、それらは単に型\code{a}の値をカプセル化するだけだ。つまり記号的に$a + 0 = a$だ。

和型はHaskellでは非常に一般的だが、C++での対応物である\code{union}や\code{variant}は
はるかに一般的でない。それにはいくつかの理由がある。

まず最も単純な和型は単なる列挙であり、C++では\code{enum}を使って実装される。Haskellの和型：

\src{snippet24}
に相当するC++は：

\begin{snip}{cpp}
enum { Red, Green, Blue };
\end{snip}
さらに単純な和型：

\src{snippet25}
はC++のプリミティブ型\code{bool}だ。

値の存在または不在をエンコードする単純な和型は、C++では空文字列、負の数、ヌルポインタなどの
特別なトリックや「不可能な」値を使って様々な方法で実装される。このような省略可能性が意図的な
場合、Haskellでは\code{Maybe}型を使って表現される：

\src{snippet26}
\code{Maybe}型は2つの型の和だ。2つのコンストラクタを個別の型に分けるとわかる。最初のものは
次のようになる：

\src{snippet27}
これは\code{Nothing}という1つの値を持つ列挙だ。言い換えれば、ユニット型\code{()}と等価な
シングルトンだ。2番目の部分：

\src{snippet28}
は単に型\code{a}のカプセル化だ。\code{Maybe}は次のようにエンコードすることもできた：

\src{snippet29}
より複雑な和型はC++ではポインタを使って偽装されることが多い。ポインタはヌルか、特定の型の
値を指すかのどちらかだ。例えば、Haskellのリスト型は（再帰的な）和型として定義できる：

\src{snippet30}
これはC++で、空リストを実装するためにヌルポインタのトリックを使って変換できる：

\begin{snip}{cpp}
template<typename A>
class List {
    Node<A> * _head;
public:
    List() : _head(nullptr) {} // Nil
    List(A a, List<A> l)       // Cons
      : _head(new Node<A>(a, l))
    {}
};
\end{snip}
Haskellの2つのコンストラクタ\code{Nil}と\code{Cons}が、類似の引数を持つ2つのオーバーロードされた
\code{List}コンストラクタに変換されることに注目しよう（\code{Nil}には引数なし；\code{Cons}には
値とリスト）。\code{List}クラスは和型の2つのコンポーネントを区別するタグを必要としない。代わりに
\code{\_head}に対して特別な\code{nullptr}値を使って\code{Nil}をエンコードする。

しかし、HaskellとC++の型の主な違いは、Haskellのデータ構造が不変であることだ。特定のコンストラクタを
使ってオブジェクトを作ると、そのオブジェクトはどのコンストラクタが使われ何の引数が渡されたかを
永遠に記憶する。よって、\code{Just "energy"}として作られた\code{Maybe}オブジェクトは
\code{Nothing}に変わることは決してない。同様に、空リストは永遠に空であり続け、3要素のリストは
常に同じ3要素を持つ。

この不変性こそが構築を可逆にするものだ。オブジェクトを与えられると、構築に使われた部品に常に
分解できる。この分解はパターンマッチングで行われ、コンストラクタをパターンとして再利用する。
コンストラクタの引数は（もしあれば）変数（または他のパターン）に置き換えられる。

\code{List}データ型は2つのコンストラクタを持つため、任意の\code{List}の分解はそれらのコンストラクタに
対応する2つのパターンを使う。1つは空の\code{Nil}リストにマッチし、もう1つは\code{Cons}で
構築されたリストにマッチする。例えば、\code{List}上の単純な関数の定義を示す：

\src{snippet31}
\code{maybeTail}の定義の最初の部分は、\code{Nil}コンストラクタをパターンとして使い
\code{Nothing}を返す。2番目の部分は\code{Cons}コンストラクタをパターンとして使う。
最初のコンストラクタ引数には興味がないためワイルドカードで置き換える。\code{Cons}の
2番目の引数は変数\code{t}に束縛される（厳密に言えば決して変化しないが、これらを変数と呼ぶことにする：
一度式に束縛されると、変数は変わらない）。戻り値は\code{Just t}だ。\code{List}がどのように作られたかに
応じて、いずれかの節にマッチする。\code{Cons}を使って作られた場合、渡された2つの引数が取得される
（最初のものは破棄される）。

さらに精巧な和型はC++でポリモーフィッククラス階層を使って実装される。共通の祖先を持つクラスの
ファミリーは、vtableが隠しタグとして機能するバリアント型の一種として理解できる。Haskellでは
コンストラクタのパターンマッチングと特化されたコードの呼び出しで行うことが、C++ではvtableポインタに
基づく仮想関数呼び出しのディスパッチによって実現される。

\code{union}を和型として使うことはC++ではほとんどない。それはunionに何を入れられるかについて
厳しい制限があるためだ。コピーコンストラクタを持つため、unionに\code{std::string}を入れることすら
できない。

\section{型の代数}

直積型と和型を別々に使えば、様々な有用なデータ構造を定義できるが、真の力は2つを組み合わせることから
来る。ここでも合成の力を呼び起こす。

これまでに発見したことをまとめよう。型システムの基礎にある2つの可換モノイド構造を見てきた：
\code{Void}を中立元とする和型と、ユニット型\code{()}を中立元とする直積型だ。これらを加算と
乗算のアナロジーとして考えたい。このアナロジーでは、\code{Void}はゼロに、ユニット\code{()}は
1に対応する。

このアナロジーをどこまで拡張できるか見てみよう。例えば、ゼロとの乗算はゼロを与えるか？
言い換えれば、1つのコンポーネントが\code{Void}である直積型は\code{Void}と同型か？
例えば、\code{Int}と\code{Void}のペアを作れるだろうか？

ペアを作るには2つの値が必要だ。整数は簡単に用意できるが、\code{Void}型の値は存在しない。
よって、任意の型\code{a}に対して、型\code{(a, Void)}は要素を持たない――値がない――ため、
\code{Void}と等価だ。言い換えれば、$a \times 0 = 0$だ。

加算と乗算を結びつけるもう一つの性質は分配法則だ：
\[a \times (b + c) = a \times b + a \times c\]
これは直積型と和型でも成り立つか？はい、成り立つ――いつものように同型を除いて。
左辺は型に対応する：

\src{snippet32}
そして右辺は型に対応する：

\src{snippet33}
一方向に変換する関数を示す：

\src{snippet34}
そしてもう一方向に変換する関数：

\src{snippet35}
\code{case of}節は関数内のパターンマッチングに使われる。各パターンの後に矢印と、パターンが
マッチしたときに評価される式が続く。例えば、値：

\src{snippet36}
で\code{prodToSum}を呼び出すと、\code{case e of}の\code{e}は\code{Left "Hi!"}に等しくなる。
それはパターン\code{Left y}にマッチし、\code{y}に\code{"Hi!"}を代入する。\code{x}はすでに
\code{2}にマッチしているため、\code{case of}節と関数全体の結果は期待通り\code{Left (2, "Hi!")}
になる。

これら2つの関数が互いの逆関数であることは証明しないが、考えてみれば当然だ！それらは単に2つの
データ構造の内容を自明に再パッケージするだけだ。同じデータで、形式が異なるだけだ。

数学者はこのように絡み合った2つのモノイドに名前をつけた：\newterm{半環}と呼ばれる。
型の減算を定義できないため、完全な\newterm{環}ではない。そのため半環は\newterm{rig}と
呼ばれることもある。これは「\emph{n}（negative）のない環（ring）」という語呂合わせだ。
それを除けば、rigを形成する自然数についての命題を型についての命題に翻訳することで多くの
ことが得られる。いくつかの興味深いエントリを含む変換表を示す：

\begin{longtable}[]{@{}ll@{}}
  \toprule
  数値         & 型\tabularnewline
  \midrule
  \endhead
  $0$          & \code{Void}\tabularnewline
  $1$          & \code{()}\tabularnewline
  $a + b$      & \code{data Either a b = Left a | Right b}\tabularnewline
  $a \times b$ & \code{(a, b)} または \code{data Pair a b = Pair a b}\tabularnewline
  $2 = 1 + 1$  & \code{data Bool = True | False}\tabularnewline
  $1 + a$      & \code{data Maybe a = Nothing | Just a}\tabularnewline
  \bottomrule
\end{longtable}

\noindent
リスト型はかなり興味深い。なぜなら方程式の解として定義されるからだ。定義しようとしている型が
方程式の両辺に現れる：

\src{snippet37}
通常の置き換えを行い、\code{List a}を\code{x}で置き換えると、次の方程式が得られる：

\begin{Verbatim}
  x = 1 + a * x
\end{Verbatim}
型の減算や除算ができないため、従来の代数的方法では解けない。しかし、右辺の\code{x}を
\code{(1 + a*x)}で置き換え続け、分配法則を使う一連の置き換えを試みることができる。
これにより次の系列が得られる：

\begin{Verbatim}
  x = 1 + a*x
  x = 1 + a*(1 + a*x) = 1 + a + a*a*x
  x = 1 + a + a*a*(1 + a*x) = 1 + a + a*a + a*a*a*x
  ...
  x = 1 + a + a*a + a*a*a + a*a*a*a...
\end{Verbatim}
最終的に無限の直積（タプル）の和になる。これは次のように解釈できる：リストは
空\code{1}か、シングルトン\code{a}か、ペア\code{a*a}か、3つ組\code{a*a*a}か、
などだ\ldots{}。これはまさにリストそのものだ――\code{a}の並びだ！

リストにはそれ以上のものがあり、ファンクターと不動点について学んだ後で、リストや他の
再帰的データ構造に戻ってくる。

記号変数を使って方程式を解く――それが代数だ！これがこれらの型に名前を与えるものだ：
代数的データ型だ。

最後に、型の代数の非常に重要な解釈の一つに触れておこう。型\code{a}と\code{b}の直積は
型\code{a}の値\emph{と}型\code{b}の値の両方を含まなければならないことに注目しよう。
つまり、両方の型が要素を持たなければならない。一方、2つの型の和は型\code{a}の値\emph{または}
型\code{b}の値を含むため、どちらか一方が要素を持てば十分だ。論理的な\emph{and}と\emph{or}も
半環を形成し、それも型理論に写すことができる：

\begin{longtable}[]{@{}ll@{}}
  \toprule
  論理                 & 型\tabularnewline
  \midrule
  \endhead
  $\mathit{false}$     & \code{Void}\tabularnewline
  $\mathit{true}$      & \code{()}\tabularnewline
  $a \mathbin{||} b$   & \code{data Either a b = Left a | Right b}\tabularnewline
  $a \mathbin{\&\&} b$ & \code{(a, b)}\tabularnewline
  \bottomrule
\end{longtable}

\noindent
このアナロジーはより深く、論理と型理論の間のカリー＝ハワード同型の基礎となっている。
関数型について話すときにこれに戻ってくる。

\section{演習}

\begin{enumerate}
  \tightlist
  \item
        \code{Maybe a}と\code{Either () a}の間の同型を示せ。
  \item
        Haskellで定義された次の和型を示す：

        \begin{snip}{haskell}
data Shape = Circle Float
           | Rect Float Float
\end{snip}
        \code{Shape}に作用する\code{area}のような関数を定義するには、2つのコンストラクタで
        パターンマッチングを行う：

        \begin{snip}{haskell}
area :: Shape -> Float
area (Circle r) = pi * r * r
area (Rect d h) = d * h
\end{snip}
        C++またはJavaでインターフェースとして\code{Shape}を実装し、\code{Circle}と
        \code{Rect}の2つのクラスを作れ。\code{area}を仮想関数として実装せよ。
  \item
        前の例を続けて：\code{Shape}の周長を計算する新しい関数\code{circ}を簡単に追加できる。
        \code{Shape}の定義に触れることなく実装できる：

        \begin{snip}{haskell}
circ :: Shape -> Float
circ (Circle r) = 2.0 * pi * r
circ (Rect d h) = 2.0 * (d + h)
\end{snip}
        C++またはJavaの実装に\code{circ}を追加せよ。元のコードのどの部分を変更する必要があったか？
  \item
        さらに続けて：\code{Shape}に新しい形状\code{Square}を追加し、必要な更新を行え。
        HaskellとC++またはJavaでどのコードに触れる必要があったか？（Haskellプログラマーでなくても、
        変更はかなり明白なはずだ。）
  \item
        $a + a = 2 \times a$が型に対して（同型を除いて）成り立つことを示せ。変換表によると、
        $2$は\code{Bool}に対応することを覚えておこう。
\end{enumerate}
